\section{Related Work}
\label{sec:related}

\para{Self-refinement at inference time.}
\citet{madaan2023selfrefine} introduced the \selfrefine framework, in which a single LLM iteratively generates, critiques, and revises its output without additional training.
They report roughly 20\% average improvement across seven tasks, though their analysis of creative writing is limited.
\citet{huang2022selfimprove} showed that LLMs can self-improve on reasoning tasks by fine-tuning on their own high-confidence outputs.
More recently, \citet{zhang2023contemplation} proposed an unsupervised RL framework where the model scores its own outputs and optimizes against those self-scores, achieving gains on both reasoning and open-ended generation.
Our work focuses specifically on narrative fiction, measuring self-refinement depth and contrasting it with self-selection strategies.

\para{Self-rewarding and self-training.}
\citet{yuan2024selfrewarding} demonstrated that a single LLM can simultaneously serve as instruction-follower and reward model, with both capabilities improving through iterative DPO training.
Their strongest gains appeared in the writing category of MT-Bench, though response length inflated by 2.3$\times$---a potential sign of reward hacking.
\citet{jia2025writingzero} addressed this for Chinese writing by introducing a generative reward model that produces pairwise critiques rather than scalar scores, reducing reward hacking by 7$\times$.
\citet{chen2024selfrewarding} further showed that self-evolved reward learning can bootstrap reward model quality alongside the generator.
Unlike these training-time approaches, we study inference-time self-critique, where the model does not permanently learn from its critiques.

\para{Critique for creative writing.}
\citet{bae2024critics} proposed \textsc{CritiCS}, a multi-agent framework with diverse critic personas for story generation.
They found that single-agent critique was insufficient and that multiple specialized critics were essential for meaningful creative improvement.
\citet{wu2025prefine} applied critique-and-refine specifically to personalized story generation with simulated user agents.
\citet{chen2025evolvr} developed a self-evolving pairwise reasoning approach for story evaluation, achieving state-of-the-art results on multiple story benchmarks and demonstrating that improved evaluation translates to improved generation when used as a reward model.
Our work differs in that we isolate single-model self-critique without multi-agent setups or training-time adaptation, directly measuring what one round of honest self-reflection achieves.

\para{LLM-as-judge and evaluation bias.}
Using LLMs to evaluate LLM outputs has become widespread~\citep{zheng2023judging}, but introduces known biases: position bias, verbosity bias, and self-enhancement bias.
\citet{pan2024rewardhacking} demonstrated that LLM evaluator scores can diverge from human judgments over refinement iterations, a phenomenon they call spontaneous reward hacking.
We mitigate these concerns through position randomization, separate evaluation contexts, and explicit instructions against favoring longer stories, though we acknowledge that shared model biases between generator and judge remain a limitation.

\para{Constitutional AI and RLAIF.}
\citet{bai2022constitutional} established the paradigm of using AI-generated feedback---including self-critique---as a substitute for human feedback in training loops.
\citet{lee2023rlaif} systematically compared RLAIF with RLHF and found comparable performance, validating AI feedback as a viable reward signal.
These works provide the theoretical foundation for self-critique as a quality signal, which we test in the specific domain of narrative fiction.
